% Ejercicio 12, inciso c) — issue #8. Mismo contenido que inciso-c.md.
% Fragmento sin preámbulo: lo incluye latex/main.tex, que ya pone el título del inciso.
% \sloppy dentro de un grupo: los nombres largos en \texttt no se salen del margen,
% y el ajuste no afecta a los incisos siguientes.
\begingroup\sloppy

\begin{quote}
\itshape
Integra con Composite para manejar grupos de cuentas (ej. «Cuentas de una sucursal») y con
Iterator para recorrer las transacciones de una cuenta. ¿Cómo afecta el Composite a la
ejecución de operaciones sobre grupos?
\end{quote}

Se integra sobre el código del inciso~b. El grupo de cuentas entra por la interfaz
\texttt{ICuenta} del inciso~a, que solo declara \texttt{Aceptar}.

\subsubsection{Composite: grupos de cuentas}

Composite compone objetos en estructuras de árbol y permite que el cliente trate igual a un
objeto individual y a una composición de objetos (Gamma et al., 1994, p.~183). Aquí el objeto
individual es una cuenta y la composición es un grupo: una sucursal, una región o el banco
entero.

\begin{table}[H]
\centering
\small
\begin{tabularx}{\textwidth}{>{\raggedright\arraybackslash}p{2.1cm}
                             >{\raggedright\arraybackslash}p{4.6cm}
                             >{\raggedright\arraybackslash}X}
\toprule
\textbf{Rol (GoF)} & \textbf{Clase} & \textbf{Qué hace} \\
\midrule
Component & \texttt{ICuenta} & Lo que tienen en común una cuenta y un grupo:
  \texttt{Aceptar(IVisitanteCuenta)}. \\
\addlinespace
Leaf & las 4 cuentas (a través de \texttt{Cuenta}) & Llaman al \texttt{Visitar} de su clase. \\
\addlinespace
Composite & \texttt{GrupoCuentas} & Guarda sus hijos (cuentas u otros grupos) y en su
  \texttt{Aceptar} le pasa el visitante a cada uno. \\
\addlinespace
Client & \texttt{ProcesadorTransacciones}, \texttt{Program.cs} & Le pasan un visitante a una
  cuenta o a un grupo sin distinguirlos. \\
\bottomrule
\end{tabularx}
\caption{Roles de Composite en la solución.}
\end{table}

\begin{lstlisting}[basicstyle=\ttfamily\footnotesize, columns=fullflexible, keepspaces=true, caption={\texttt{GrupoCuentas.Aceptar}: el recorrido recursivo del árbol.}]
public void Aceptar(IVisitanteCuenta visitante)
{
    visitante.EntrarEnGrupo(this);
    foreach (ICuenta hijo in _hijos)
    {
        hijo.Aceptar(visitante);   // si el hijo es un grupo, recorre a su vez los suyos
    }

    visitante.SalirDeGrupo(this);
}
\end{lstlisting}

\begin{itemize}
  \item \textbf{El cliente trata igual a una cuenta y a un grupo.} \texttt{cuenta.Aceptar(v)}
        y \texttt{banco.Aceptar(v)} se escriben igual, y
        \texttt{EjecutarEnCadaCuenta(\allowbreak{}ICuenta destino, \ldots)} acepta cualquiera de los dos.
        GoF lo presenta como la ventaja principal del patrón: el cliente no sabe, ni le
        importa, si tiene una hoja o un compuesto (p.~186).
  \item \textbf{Los grupos se anidan} (sucursal $\rightarrow$ región $\rightarrow$ banco), y
        el recorrido es \textbf{recursivo}: el \texttt{Aceptar} de un grupo llama al
        \texttt{Aceptar} de cada hijo. \texttt{Agregar} rechaza los ciclos (un grupo dentro de
        sí mismo o de uno de sus descendientes), porque el recorrido no terminaría.
  \item \textbf{Es el uso de Visitor que GoF señala} (p.~380): el elemento compuesto recorre
        sus hijos en su \texttt{Accept} (p.~373), y así el recorrido del árbol es
        responsabilidad de la estructura, no del visitante (p.~374).
\end{itemize}

\paragraph{Cuánto costó agregar el Composite.}
\texttt{IVisitanteCuenta} ganó dos métodos, \texttt{EntrarEnGrupo} y \texttt{SalirDeGrupo},
con una \textbf{implementación vacía por omisión}. Los visitantes que no distinguen grupos
(ejecutar, revertir, auditar y calcular intereses) heredan esa implementación y \textbf{no
cambiaron}. GoF admite esta implementación por omisión como excepción (p.~371). El costo
total fue:
\begin{itemize}
  \item 1 clase nueva (\texttt{GrupoCuentas});
  \item 1 interfaz con 2 métodos nuevos;
  \item 0 cuentas y 0 visitantes modificados.
\end{itemize}
El inciso~d usa estos números.

\subsubsection{Iterator: las transacciones de una cuenta}

Iterator permite recorrer los elementos de un agregado \textbf{sin exponer su representación
interna} (Gamma et al., 1994, p.~289), y admitir varios recorridos distintos del mismo
agregado (p.~291). En C\# es \texttt{IEnumerable<T>}, como en las colecciones de .NET que
muestra la conferencia (diap.~35--37).

\begin{table}[H]
\centering
\small
\begin{tabularx}{\textwidth}{>{\raggedright\arraybackslash}p{2.6cm} >{\raggedright\arraybackslash}X}
\toprule
\textbf{Rol (GoF)} & \textbf{Aquí} \\
\midrule
Aggregate & \texttt{Cuenta}, que guarda su historial de asientos en una lista privada \\
CreateIterator & \texttt{Historial()}, y también \texttt{GrupoCuentas.Hijos} para los hijos de
                 un grupo \\
Iterator & el \texttt{IEnumerator<Asiento>} que genera el compilador a partir de
           \texttt{yield return} \\
\bottomrule
\end{tabularx}
\caption{Roles de Iterator en la solución.}
\end{table}

\begin{lstlisting}[basicstyle=\ttfamily\footnotesize, columns=fullflexible, keepspaces=true]
public IEnumerable<Asiento> Historial()
{
    foreach (Asiento asiento in _asientos)
    {
        yield return asiento;
    }
}
\end{lstlisting}

\begin{itemize}
  \item \textbf{No se expone la lista interna.} \texttt{Historial()} no devuelve la
        \texttt{List<Asiento>}: convertirla a \texttt{ICollection<Asiento>} para agregar o
        quitar asientos da \texttt{null} (hay una prueba).
  \item \textbf{Hay iteradores filtrados:} \texttt{Historial(\allowbreak{}desde, hasta)} filtra por fecha,
        y \texttt{Historial<\allowbreak{}Transferencia>()} por tipo de transacción.
  \item \textbf{Incluye los asientos revertidos}, marcados como tales: el historial es el
        rastro completo de la cuenta.
  \item \textbf{Si la cuenta cambia mientras se recorre, el recorrido falla} en vez de
        devolver datos a medias. Esto importa para la restricción (\#10), donde varias
        transacciones corren en paralelo.
  \item \textbf{Interno y externo.} \texttt{Historial()} es un iterador \emph{externo}: el
        cliente controla el bucle. El recorrido del árbol de grupos, en cambio, es
        \emph{interno}: lo hace \texttt{Aceptar}. GoF explica que los iteradores externos son
        difíciles de implementar sobre estructuras recursivas como las de Composite, porque la
        posición abarca varios niveles (p.~295), y por eso el árbol lo recorre el propio
        Composite.
\end{itemize}

\subsubsection{¿Cómo afecta el Composite a la ejecución de operaciones sobre grupos?}

\paragraph{1. La operación llega a todas las hojas, en profundidad y en el orden de inserción.}
Una operación sobre un grupo se aplica a \textbf{todas las cuentas del árbol}, no a los
grupos: los grupos solo reenvían el visitante. El orden es un recorrido en profundidad que
respeta el orden en que se agregaron los hijos. Con el árbol de la simulación:

\begin{lstlisting}[language={}, numbers=none, basicstyle=\ttfamily\footnotesize, columns=fullflexible, keepspaces=true]
Banco
+- Region Occidente
|  +- Sucursal Centro:  AHO-001, COR-001
|  +- Sucursal Playa:   INV-001, AHO-002
+- Sucursal Empresas:   EMP-001
\end{lstlisting}

el orden es AHO-001, COR-001, INV-001, AHO-002 y EMP-001. El visitante recibe además
\texttt{EntrarEnGrupo} y \texttt{SalirDeGrupo} alrededor de cada grupo; la prueba
\texttt{Aceptar\_\allowbreak{}RecorreLosGruposAnidados\allowbreak{}EnProfundidadYEnOrden} fija esa secuencia completa.

\paragraph{2. Si falla en una cuenta del grupo: mejor esfuerzo con reporte.}
\texttt{ProcesadorTransacciones.\allowbreak{}EjecutarEnCadaCuenta(\allowbreak{}destino, crear)} aplica una operación a
una cuenta o a un grupo. Funciona así:
\begin{enumerate}
  \item Obtiene las cuentas del destino con \texttt{RecolectarCuentas}, un visitante que
        recorre el árbol.
  \item Crea \textbf{una transacción por cuenta} y la ejecuta con el «todo o nada» del
        inciso~b.
  \item Si una cuenta rechaza la suya, \textbf{las demás siguen}, y el resultado de cada una
        (aplicada, o rechazada con su motivo) queda en el reporte.
\end{enumerate}

\textbf{Por qué mejor esfuerzo y no todo o nada para el grupo:}
\begin{itemize}
  \item \textbf{Cada cuenta es de un cliente distinto.} Un cobro de mantenimiento a una
        sucursal son N hechos independientes, no uno solo. Si un cliente no tiene fondos, no
        hay razón para dejar de cobrarles a los demás.
  \item \textbf{Con todo o nada, una sola cuenta bloquearía al banco entero.} En la
        simulación, 2 de 5 cuentas rechazan el cobro (una agotó sus débitos del mes y otra no
        tiene saldo). Con todo o nada no se habría cobrado a ninguna, y con miles de cuentas
        casi siempre habrá alguna que rechace.
  \item \textbf{No se pierde ni se oculta nada.} Cada transacción por cuenta sigue siendo
        atómica, una cuenta que rechaza queda intacta, y el reporte dice qué cuentas fallaron
        y por qué, para reintentarlas o avisar al cliente.
\end{itemize}

\textbf{Cuándo sí conviene todo o nada:} cuando la operación es \textbf{un solo hecho} que
toca varias cuentas, como una transferencia o una nómina (un débito en la empresa y un crédito
a cada empleado). Ahí no tiene sentido aplicar una parte. Eso se modela como \textbf{una}
transacción con varios movimientos, y \texttt{ProcesadorTransacciones.Ejecutar} ya la aplica
con todo o nada (inciso~b).

\begin{table}[H]
\centering
\small
\begin{tabularx}{\textwidth}{>{\raggedright\arraybackslash}p{2.6cm}
                             >{\raggedright\arraybackslash}X
                             >{\raggedright\arraybackslash}X}
\toprule
 & \textbf{Mejor esfuerzo (\texttt{EjecutarEnCadaCuenta})} & \textbf{Todo o nada
   (\texttt{Ejecutar})} \\
\midrule
Qué es & N transacciones independientes, una por cuenta & 1 transacción con N movimientos \\
\addlinespace
Si una cuenta rechaza & las demás se aplican; el reporte la señala & se revierten las ya
  aplicadas \\
\addlinespace
Ejemplos & cobro de mantenimiento, abono de intereses & transferencia, pago de nómina \\
\bottomrule
\end{tabularx}
\caption{Las dos políticas ante un fallo en una cuenta.}
\end{table}

\textbf{Por qué no basta con \texttt{grupo.Aceptar(\allowbreak{}new EjecutarTransaccion(t))}:}
\texttt{EjecutarTransaccion} aplica los movimientos de \emph{una} transacción, que no tiene
movimientos en las demás cuentas del grupo. Además, el procesador necesita el resultado de
cada cuenta, para reportarlo o para revertirlo. Por eso primero obtiene las hojas con un
visitante y después trabaja cuenta por cuenta.

\paragraph{3. Cómo un visitante acumula resultados al recorrer un grupo.}
\texttt{ResumirPorGrupo} calcula en \textbf{un solo recorrido} los totales de cada grupo:
cuántas cuentas tiene, su saldo y los cargos cobrados.
\begin{itemize}
  \item En \texttt{EntrarEnGrupo} abre un acumulador para el grupo, y en
        \texttt{SalirDeGrupo} lo cierra.
  \item Cada cuenta visitada \textbf{suma en todos los grupos abiertos}, es decir, en todos
        los que la contienen: una cuenta de la Sucursal Centro cuenta en su sucursal, en la
        Región Occidente y en el Banco.
  \item \textbf{Cada tipo de cuenta aporta sus cargos a su manera:} ahorro no cobra;
        corriente aporta el cargo por cada vez que quedó en sobregiro; inversión, sus
        penalizaciones; empresarial, sus comisiones. Es el doble despacho otra vez.
\end{itemize}

Es la consecuencia de Visitor que GoF llama acumular estado durante el recorrido (p.~372). Un
visitante que no distingue grupos también acumula: \texttt{CalcularIntereses}, del inciso~b,
calcula los intereses de todas las cuentas del banco con \texttt{banco.Aceptar(\ldots)},
\textbf{sin ningún cambio}.

\subsubsection{Resultados de la simulación}

\texttt{Program.cs} agrega las secciones 5 y 6 al final de la simulación del inciso~b, con una
cuenta más (AHO-002, con 3.00 de saldo):

\begin{lstlisting}[language={}, numbers=none, basicstyle=\ttfamily\scriptsize, columns=fullflexible, keepspaces=true]
--- 5. Composite: totales por grupo en un solo recorrido (ResumirPorGrupo) ---
  Banco                  5 cuentas  saldo  69022.00  cargos cobrados 101.00
    Región Occidente     4 cuentas  saldo  35303.00  cargos cobrados  20.00
      Sucursal Centro    2 cuentas  saldo   1320.00  cargos cobrados   0.00
      Sucursal Playa     2 cuentas  saldo  33983.00  cargos cobrados  20.00
    Sucursal Empresas    1 cuenta   saldo  33719.00  cargos cobrados  81.00

  CalcularIntereses sobre el banco entero (mismo visitante, sin cambios): total 183.89

  Cobro de mantenimiento de 5.00 a todas las cuentas del banco (mejor esfuerzo):
  Rechazada  Pago de Mantenimiento por 5.00 desde AHO-001
             motivo: AHO-001: alcanzó el límite de 3 retiros del mes
  Aplicada   Pago de Mantenimiento por 5.00 desde COR-001
  Aplicada   Pago de Mantenimiento por 5.00 desde INV-001
  Rechazada  Pago de Mantenimiento por 5.00 desde AHO-002
             motivo: AHO-002: saldo insuficiente (saldo 3.00, a debitar 5.00)
  Aplicada   Pago de Mantenimiento por 5.00 desde EMP-001
  → 3 aplicadas y 2 rechazadas; las rechazadas no tocaron sus cuentas.

--- 6. Iterator: historial de COR-001 ---
  Retiro de 600.00 de COR-001                    Debito  efecto  -605.00  (revertido)
  Transferencia de 50.00 de COR-001 a INV-001    Debito  efecto   -55.00  (revertido)
  Pago de Mantenimiento por 5.00 desde COR-001   Debito  efecto    -5.00

  Solo las transferencias de COR-001:
  Transferencia de 50.00 de COR-001 a INV-001
\end{lstlisting}

\begin{itemize}
  \item \textbf{Los totales cuadran por niveles:} Centro (1320.00) más Playa (33983.00) dan
        la Región (35303.00), y la Región más Empresas (33719.00) dan el Banco (69022.00).
  \item \textbf{El cobro sigue la política del punto 2:} las 3 cuentas que pueden pagar
        pagan; AHO-001 (ya hizo 3 débitos este mes) y AHO-002 (sin saldo) quedan intactas y
        aparecen en el reporte con su motivo.
  \item \textbf{El historial de COR-001} muestra los dos débitos revertidos del inciso~b y el
        cobro de mantenimiento; el filtro por tipo deja solo la transferencia.
\end{itemize}

\subsubsection{Pruebas}

\texttt{dotnet test} ejecuta \textbf{57 pruebas}, todas correctas: las 46 del inciso~b y 11
nuevas.

\begin{table}[H]
\centering
\footnotesize
\begin{tabularx}{\textwidth}{>{\raggedright\arraybackslash}p{7.2cm} >{\raggedright\arraybackslash}X}
\toprule
\textbf{Prueba} & \textbf{Qué comprueba} \\
\midrule
\texttt{Aceptar\_RecorreLosGruposAnidados\-EnProfundidadYEnOrden} & Un grupo anidado en tres
  niveles recorre todas sus hojas en profundidad y en orden, con los avisos de entrada y
  salida. \\
\texttt{Agregar\_RechazaLosCiclos} & No se puede meter un grupo dentro de sí mismo ni dentro de
  un descendiente. \\
\texttt{Hijos\_NoExponeLaListaInterna} & Los hijos se recorren sin poder modificar la lista. \\
\texttt{ResumirPorGrupo\_SumaCadaCuentaEn\-TodosLosGruposQueLaContienen} & Totales de saldo y
  cargos correctos en cada nivel de un árbol anidado. \\
\texttt{UnVisitanteQueNoDistingueGrupos\_\-FuncionaSobreUnGrupoSinCambios} &
  \texttt{CalcularIntereses} recorre un grupo entero sin ningún cambio. \\
\texttt{EjecutarEnCadaCuenta\_ConFalloParcial\-SigueConLasDemasYReportaCadaResultado} &
  \textbf{Fallo parcial:} la cuenta que rechaza queda intacta, las demás se cobran y el reporte
  da cada resultado. \\
\texttt{EjecutarEnCadaCuenta\_TrataIgual\-AUnaCuentaSola} & La misma operación funciona con una
  cuenta sola. \\
\texttt{Historial\_RecorreLosAsientosEnOrden\-IncluidosLosRevertidos} & El historial sale en
  orden, incluidos los asientos revertidos. \\
\texttt{Historial\_NoExponeLaListaInterna} & El historial no se puede convertir en una
  colección modificable. \\
\texttt{Historial\_FiltraPorFecha} & El filtro por fecha devuelve solo los asientos del
  intervalo. \\
\texttt{Historial\_FiltraPorTipoDeTransaccion} & El filtro por tipo devuelve solo las
  transferencias. \\
\bottomrule
\end{tabularx}
\caption{Pruebas nuevas del inciso c.}
\end{table}

\paragraph{Referencias.}
Gamma, E., Helm, R., Johnson, R., y Vlissides, J. (1994). \emph{Design patterns: Elements of
reusable object-oriented software}. Addison-Wesley. Los números de página corresponden a la
copia del libro en la bibliografía de la asignatura (\emph{Documentacion Patrones}), igual que
en los incisos~a y~b.

Ingeniería de Software. \emph{Marcos de Trabajo y Patrones de Diseño} [diapositivas de la
conferencia].
\endgroup
